\section{An incomplete odd regular core}

\begin{theorem}[982-R07, 982-R08]
Let \(A\) consist of \(m\) vertices of a unit regular odd \(M\)-gon, and \(B\)
consist of \(t\) points off the unit circle, with \(A\cup B\) strictly convex.
If \(M\ge11\) and \(\cos h<|b|\le\cos h/\cos(2h)\) for all \(b\in B\),
where \(h=\pi/M\), \(m\ge3t/2+83/2\) guarantees a qualifying pin in \(A\).
If \(B\) is cocircular, then \(M\ge3\) and \(m\ge t+44\) suffice without the annular hypothesis.
\end{theorem}

\begin{proof}
First, any three distinct lines spanned by vertices of an odd regular polygon
that meet at a finite point meet at a polygon vertex. For concurrence at
a nonvertex the endpoints are all distinct. Normalize the three pairs to
\((1,r),(a,b),(c,d)\) with entries in \(\mu_M\). The equation of the second
line is \(z+ab\bar z=a+b\). Substituting \(z=1+\lambda(r-1)\) with real
\(\lambda\), and comparing the third line, gives
\[
 (1-a)(1-b)(1-cd/r)=(1-c)(1-d)(1-ab/r)\ne0.
\]
The nonzero assertion uses nonparallelism and distinct endpoints.
Apply the cyclotomic field automorphism squaring the roots, then divide
by this equality. The resulting equality of products with plus signs,
together with the common product \(abcd/r\), forces equality of the
first two elementary symmetric coefficients. Thus \((a,b,cd/r)\)
and \((c,d,ab/r)\) are the same multiset. Since \(a\notin\{c,d\}\) this
forces \(b=r\), a contradiction.

For the annular case put \(S=\sin h,\ C=\cos h,\ U=C/\cos(2h)\).
Form a graph on \(B\) by joining pairs whose bisector meets two vertices
of \(A\). A triangle of this graph would have its circumcenter at
a vertex of \(A\), by the preceding concurrence property. Rotate this
common vertex to 1 and write the three targets as \(1+R e^{i\theta_i}\).
They lie in one open half-plane at 1, as do all other vertices, so
choose \(\theta_1<\theta_2<\theta_3<\theta_1+\pi\). The second unit-circle
intersection of the \(i,j\) bisector is
\(-e^{i(\theta_i+\theta_j)}\in A\), lying on the positive ray at angle
\((\theta_i+\theta_j)/2\) from 1. For pairs \(12,23\) these rays enclose
\(\theta_2\), with angular gap less than \(\pi/2\). Both unit-circle
intersections are at least \(2S\) from 1. Their joining segment has
projection strictly greater than \(\sqrt 2 S\) along the mid-direction.
Since the middle target must be outside the triangle bounded by this
segment and the vertex 1, \(R>\sqrt2 S\).

From the three second intersections being roots of unity, the
\(\theta_i\) are congruent modulo \(2h\) and their pairwise sums are
\(\pi\) modulo \(2h\). The folded angles
\(\gamma_i=\arccos(\cos\theta_i)\) are distinct half-integer multiples
of \(h\) in \([h/2,\pi-h/2]\). In fact two original angles giving
the same folded angle would either coincide or sum to zero modulo
\(2\pi\), incompatible with the pairwise sum congruence for odd \(M\).
Thus the maximum and minimum folded angles differ by at least \(2h\).

The allowed squared-radius width is
\[
 U^2-C^2=4S^2\left(\frac{1-S^2}{1-2S^2}\right)^2<\frac{121}{25}S^2
\]
since \(h\le\pi/11<2/7\) gives \(S^2<1/12\). The actual width equals
\(2R\) times the span of the cosines of \(\gamma_i\). For \(R\ge1\)
the least span over any interval of length \(2h\) in
\([h/2,\pi-h/2]\) is \(2\sin(3h/2)S\), and
\(\sin(3h/2)/S>3\sqrt3/4>121/100\). This already exceeds the allowed width.
For \(\sqrt2 S<R\le1\), the radial bounds give
\[
 -\frac{R^2+S^2}{2R}<\cos\gamma_i\le\frac{U^2-1}{2R}<\frac{54S^2}{25R}.
\]
Here \(U^2-1=(3S^2-4S^4)/(1-2S^2)^2\).
If \(R\ge3S/2\), both extreme cosine magnitudes are bounded
by \(13/24\) using \(S^2<1/12\) and monotonicity of \(R+S^2/R\).
Then the sine between the extreme angles exceeds \(5/6\), so their
squared-radius span exceeds \(5Sh>121S^2/25\). If \(R\le3S/2\),
both magnitudes are below \(1/2\) using \(\sqrt2 S<R\), so the width
exceeds \(2\sqrt6 Sh>121S^2/25\). This contradicts the annulus in
all cases. The graph is triangle-free, hence has at most
\(\lfloor t^2/4\rfloor\) edges by Mantel's inequality.

Each nonzero \(b\in B\) causes at most 22 merges with old circle
distances over pins in \(A\). The five-term argument for the majority
theorem bounds the number of oriented torsion solutions by 44, and
each merge supplies distinct \(u,u^{-1}\): \(u=1\) would place \(b\)
on the circle, and \(u=-1\) is excluded by odd \(M\).
If the origin belongs to \(B\), strict support there confines \(A\)
to an open semicircle and an extreme pin of \(A\) suffices by the
threshold. Hence we may exclude that case. In the annular case the
total double-witness count from pairs of new targets is bounded by
\(\binom t2+\lfloor t^2/4\rfloor\), counting at most one per pair and
one extra per graph edge. Failure at all core pins would require
\[
 m(t+1)/2\le22t+\binom t2+\lfloor t^2/4\rfloor
 \le 3t^2/4+43t/2.
\]
This contradicts \(m\ge3t/2+83/2\).

For the cocircular case let \(Q\) be the center of \(B\). If \(Q\notin A\),
all bisectors of \(B\)-pairs pass through \(Q\), with at most two lines
meeting \(A\) twice by the concurrence result. If \(Q\) is itself an
unselected polygon vertex, no such line can meet \(A\) twice. Each line
accounts for at most \(\lfloor t/2\rfloor\) disjoint reflected pairs.
Consequently total pair merges are at most \(\binom t2+2\lfloor t/2\rfloor\).
With the 22-target penalty the failure condition gives
\(m(t+1)\le t^2+45t\), impossible for \(m\ge t+44\).
If \(Q\in A\), omit its pin but keep all original \(A\)-targets.
Each pair of new targets then accounts for at most one merge, since
one of its two circle intersections is \(Q\). Thus failure requires
\((m-1)(t+1)\le44t+2\binom t2=t^2+43t\), also impossible.
\end{proof}
